% Lösungen Zuordnungen (zusammenbau v0.4, Heft)
\einheitenkopf{Zuordnungen}

\erg{47}{a) (1|15), (2|30), (3|45) \quad b) (1|2), (2|4), (3|6), (4|8) \quad c) Mia hat die Achsen vertauscht: richtig sind (1|3) und (2|6), nicht (3|1) und (6|2).}
\erg{48}{a) (0|70), (1|95), (2|135), (3|185), (4|260) \quad b) (0|40), (1|60), (2|85), (3|125), (4|180)}
\erg{49}{a) 5 (20 : 4 = 5) \quad b) 1 Kästchen = 5 km, denn 60 km sind dann 12 Kästchen. \quad c) zum Beispiel 1 Kästchen = 1 Heft, Achse von 0 bis 6 beschriftet}
\erg{50}{a) fehlende Werte 90 l und 30 l; zum Beispiel 1 Kästchen = 2 min (30 min = 15 Kästchen) und 1 Kästchen = 5 l (90 l = 18 Kästchen); Punkte (0|90), (5|80), (10|70), (15|60), (20|50), (25|40), (30|30) \quad b) fehlende Werte 15 \% und 75 \%; zum Beispiel 1 Kästchen = 5 min (60 min = 12 Kästchen) und 1 Kästchen = 5 \% (75 \% = 15 Kästchen); Punkte (0|15), (10|25), (20|35), (30|45), (40|55), (50|65), (60|75)}
\erg{51}{a) weniger \quad b) 5 Tage \quad c) 3 h}
\erg{52}{a) 6 · 10 = 60 Ziegentage; 60 : 5 = 12 Tage \quad b) 3 · 10 = 30 Wandertage; 30 : 5 = 6 Tage}
\erg{53}{a) 1,79 € je Liter \quad b) 4 · 2,50 € = 10 € \quad c) 6 : 1,50 = 4 kg}
\erg{54}{a) 30\,000 : 100 · 8 = 2\,400 l; 2\,400 · 1,625 € = 3\,900 € \quad b) 15 · 400 = 6\,000 l; 6\,000 · 1,48 € = 8\,880 € \quad c) 50\,000 : 100 · 7 = 3\,500 l; 3\,500 · 0,08 € = 280 €; die Aussage stimmt. \quad d) 200 · 90 = 18\,000 l; 18\,000 · 0,10 € = 1\,800 €; die Aussage stimmt. \quad e) 25 · 0,4 ct = 10 ct = 0,10 € \quad f) 30 · 1,2 ct = 36 ct = 0,36 € \quad g) 310 − 150 = 160 m; 160 : 40 = 4 m/s; 4 · 3,6 = 14,4 km/h \quad h) 800 − 400 = 400 m; 400 : 200 = 2 m/s; 2 · 3,6 = 7,2 km/h}
\erg{55}{a) 3 km : 4 km/h = 0,75 h = 45 min; 9:50 Uhr + 30 min + 45 min = 11:05 Uhr \quad b) 9 km : 18 km/h = 0,5 h = 30 min; 14:40 Uhr + 20 min + 30 min = 15:30 Uhr \quad c) $2{,}25 \cdot 10^{8}$ km : 300\,000 km/s = 750 s; 750 s : 60 = 12,5 min \quad d) $1{,}5 \cdot 10^{8}$ km : 300\,000 km/s = 500 s; 500 s : 60 ≈ 8,3 min \quad e) 630 : 3,5 = 180 s = 3 min \quad f) 420 : 1,4 = 300 s = 5 min}
